%%%%%%%%%%%%%%%%%%%%%%%%%%%%%%%%%%%%%%%%%%%%%%%%%%%%%%%%%%
% FILE: chapter03_escribir_contenido.tex
%
% PURPOSE
% -------------------------------------------------------
% Explains how to write the main content of the document
% using the template structure.
%
%
% ARCHITECTURE ROLE
% -------------------------------------------------------
% User guide chapter.
%
% Describes the workflow authors should follow when
% writing documents with this template.
%
%
% DEPENDENCIES
% -------------------------------------------------------
% environments.tex
% listings
%
%
% NOTES FOR DEVELOPERS
% -------------------------------------------------------
% This chapter explains how users should interact with
% the template when writing content.
%
% It should not document template internals.
%
%%%%%%%%%%%%%%%%%%%%%%%%%%%%%%%%%%%%%%%%%%%%%%%%%%%%%%%%%%


\chapter{Cómo escribir contenido}

El contenido principal del documento se escribe en la
carpeta:

\texttt{02\_chapters}

Cada archivo dentro de esta carpeta representa un capítulo
independiente del documento.



% --------------------------------------------------
% ORGANIZACIÓN DE LOS CAPÍTULOS
% --------------------------------------------------

\section{Organización de los capítulos}

Un capítulo se define mediante el comando:

\begin{lstlisting}[caption={Creación de un capítulo},label={code:chapter}]
\chapter{Titulo del capitulo}
\end{lstlisting}

\source{Plantilla}

Cada capítulo debe almacenarse en un archivo independiente
dentro de la carpeta \texttt{02\_chapters}.



\begin{definitionbox}[Capítulo]
Un capítulo es la unidad principal de organización
del contenido dentro de un documento largo.
\end{definitionbox}



% --------------------------------------------------
% SECCIONES
% --------------------------------------------------

\section{Secciones}

Las secciones permiten dividir un capítulo en bloques
temáticos.

\begin{lstlisting}[caption={Creación de una sección},label={code:section}]
\section{Titulo de la seccion}
\end{lstlisting}

\source{Plantilla}



% --------------------------------------------------
% SUBSECCIONES
% --------------------------------------------------

\section{Subsecciones}

Las subsecciones permiten organizar el contenido con
mayor nivel de detalle.

\begin{lstlisting}[caption={Creación de una subsección},label={code:subsection}]
\subsection{Titulo de la subseccion}
\end{lstlisting}

\source{Plantilla}



\begin{notebox}
Se recomienda no utilizar más de tres niveles
de profundidad en la estructura del documento
para mantener la claridad del texto.
\end{notebox}



% --------------------------------------------------
% EJEMPLO DE ESTRUCTURA
% --------------------------------------------------

\section{Ejemplo de estructura}

Un capítulo típico puede organizarse de la siguiente forma.

\begin{lstlisting}[caption={Estructura típica de un capítulo},label={code:chapterstructure}]
\chapter{Introduccion}

\section{Contexto}

Texto introductorio.

\section{Objetivos}

\subsection{Objetivo general}

Descripcion.

\subsection{Objetivos especificos}

Descripcion.
\end{lstlisting}

\source{Plantilla}



% --------------------------------------------------
% FLUJO DE ESCRITURA
% --------------------------------------------------

\section{Flujo de escritura recomendado}

El flujo de trabajo recomendado al escribir un documento
con esta plantilla es el siguiente.

\begin{enumerate}

\item Crear un archivo nuevo en la carpeta
\texttt{02\_chapters}

\item Añadir el capítulo en \texttt{main.tex}

\item Escribir el contenido utilizando secciones
y subsecciones

\end{enumerate}



\begin{lstlisting}[caption={Carga de capítulos en main.tex},label={code:chapters}]
%%%%%%%%%%%%%%%%%%%%%%%%%%%%%%%%%%%%%%%%%%%%%%%%%%%%%%%%%%
% FILE: chapter01_introduccion.tex
%
% PURPOSE
% -------------------------------------------------------
% Introduces the LaTeX template and explains its
% objectives, philosophy and intended use.
%
%
% ARCHITECTURE ROLE
% -------------------------------------------------------
% First chapter of the template documentation.
%
% Provides conceptual context before explaining
% the technical structure of the project.
%
%
% DEPENDENCIES
% -------------------------------------------------------
% environments.tex
% (definitionbox, notebox, warningbox)
%
%
% NOTES FOR DEVELOPERS
% -------------------------------------------------------
% This chapter should remain conceptual.
%
% Do not include implementation details here.
% Technical explanations belong to later chapters.
%
%%%%%%%%%%%%%%%%%%%%%%%%%%%%%%%%%%%%%%%%%%%%%%%%%%%%%%%%%%


\chapter{Introducción}

Esta plantilla ha sido diseñada para facilitar la redacción de
documentos académicos y técnicos utilizando \LaTeX.

El objetivo principal es proporcionar una estructura modular
que permita separar claramente:

\begin{itemize}
\item configuración del documento
\item estilos tipográficos
\item contenido académico
\item recursos gráficos y bibliográficos
\end{itemize}

Gracias a esta separación, el autor puede concentrarse
exclusivamente en escribir el contenido del documento
sin preocuparse por el formato.



% --------------------------------------------------
% OBJETIVO DE LA PLANTILLA
% --------------------------------------------------

\section{Objetivo de la plantilla}

\begin{definitionbox}[Plantilla académica]
Una plantilla académica es una estructura predefinida
que automatiza la configuración tipográfica y
organizativa de un documento.
\end{definitionbox}

Esta plantilla está diseñada para:

\begin{itemize}
\item trabajos universitarios
\item informes técnicos
\item proyectos de ingeniería
\item Trabajos Fin de Grado (TFG)
\item Trabajos Fin de Máster (TFM)
\end{itemize}



% --------------------------------------------------
% FILOSOFÍA DE DISEÑO
% --------------------------------------------------

\section{Filosofía de diseño}

La plantilla sigue varios principios de diseño inspirados
en prácticas de ingeniería de software.

\begin{itemize}
\item modularidad
\item separación de responsabilidades
\item reutilización
\item documentación interna
\end{itemize}



\begin{notebox}
El objetivo no es únicamente generar un documento
correctamente formateado, sino también demostrar
buenas prácticas en la organización de proyectos
\LaTeX.
\end{notebox}



% --------------------------------------------------
% ORGANIZACIÓN DEL PROYECTO
% --------------------------------------------------

\section{Organización del proyecto}

El proyecto se organiza en varias carpetas que separan
la configuración del documento del contenido académico.



\begin{lstlisting}[caption={Estructura general del proyecto},label={code:estructura}]
main.tex

latex-university-template/
|-- main.tex
|-- 00_config/
|   |-- config.tex
|   `-- packages.tex
|-- 01_frontmatter/
|-- 02_chapters/
|-- 03_figures/
|-- 04_bibliography/
|-- 05_appendices/
|-- 06_styles/
`-- 07_acronyms/
\end{lstlisting}

\source{Plantilla}



% --------------------------------------------------
% USO RECOMENDADO
% --------------------------------------------------

\section{Uso recomendado}

El flujo de trabajo recomendado es el siguiente:

\begin{enumerate}
\item Configurar los metadatos del documento
\item Escribir el contenido en la carpeta \texttt{02\_chapters}
\item Añadir figuras en \texttt{03\_figures}
\item Añadir referencias en \texttt{05\_bibliography}
\end{enumerate}



\begin{warningbox}
Modificar directamente los archivos de configuración
o de estilos sin comprender su función puede producir
errores de compilación o inconsistencias tipográficas.
\end{warningbox}



% --------------------------------------------------
% METODOLOGÍA DE DOCUMENTACIÓN
% --------------------------------------------------

\section{Metodología de documentación}

Todos los archivos de la plantilla siguen una convención
de comentarios estructurados.

Cada archivo comienza con un bloque de documentación
que describe su función dentro del proyecto.

\begin{lstlisting}[caption={Bloque de documentación estándar},label={code:docblock}]
%%%%%%%%%%%%%%%%%%%%%%%%%%%%%%%%%%%%%%%%%%%%%%%%%%%%%%%%%%
% FILE: filename.tex
%
% PURPOSE
% Describe el proposito del archivo
%
% ARCHITECTURE ROLE
% Describe su funcion dentro del sistema
%
% DEPENDENCIES
% Archivos o paquetes necesarios
%
% NOTES FOR DEVELOPERS
% Indicaciones para mantenimiento
%%%%%%%%%%%%%%%%%%%%%%%%%%%%%%%%%%%%%%%%%%%%%%%%%%%%%%%%%%
\end{lstlisting}

\source{Plantilla}



\begin{examplebox}
Este sistema de documentación permite comprender
rápidamente la arquitectura del proyecto al revisar
los archivos de la plantilla.
\end{examplebox}

%%%%%%%%%%%%%%%%%%%%%%%%%%%%%%%%%%%%%%%%%%%%%%%%%%%%%%%%%%
% FILE: chapter02_arquitectura.tex
%
% PURPOSE
% -------------------------------------------------------
% Explains the internal architecture of the LaTeX template.
%
%
% ARCHITECTURE ROLE
% -------------------------------------------------------
% Technical documentation of the template structure.
%
%
% DEPENDENCIES
% -------------------------------------------------------
% environments.tex
%
%
% NOTES FOR DEVELOPERS
% -------------------------------------------------------
% If the directory structure of the template changes,
% this chapter must be updated accordingly.
%
%%%%%%%%%%%%%%%%%%%%%%%%%%%%%%%%%%%%%%%%%%%%%%%%%%%%%%%%%%


\chapter{Arquitectura del proyecto}

La plantilla está diseñada siguiendo principios de
arquitectura modular inspirados en prácticas de
ingeniería de software.

Cada componente del proyecto tiene una responsabilidad
claramente definida dentro del sistema.



% --------------------------------------------------
% PRINCIPIOS DE ARQUITECTURA
% --------------------------------------------------

\section{Principios de arquitectura}

\begin{definitionbox}[Arquitectura modular]
Una arquitectura modular divide un sistema en componentes
independientes que pueden mantenerse o modificarse de
forma aislada.
\end{definitionbox}

En esta plantilla, el proyecto se divide en varios módulos
que gestionan aspectos distintos del documento.



\begin{itemize}

\item configuración global
\item estilos tipográficos
\item contenido académico
\item recursos gráficos
\item bibliografía
\item acrónimos

\end{itemize}



\begin{notebox}
Separar estas responsabilidades permite modificar el
formato del documento sin alterar el contenido.
\end{notebox}



% --------------------------------------------------
% ESTRUCTURA DEL PROYECTO
% --------------------------------------------------

\section{Estructura del proyecto}

La estructura de carpetas del proyecto es la siguiente.

\begin{lstlisting}[caption={Estructura de directorios del proyecto},label={code:projectstructure}]
latex-university-template/
|-- main.tex
|-- 00_config/
|   |-- config.tex
|   `-- packages.tex
|-- 01_frontmatter/
|-- 02_chapters/
|-- 03_figures/
|-- 04_bibliography/
|-- 05_appendices/
|-- 06_styles/
`-- 07_acronyms/
\end{lstlisting}

\source{Plantilla}

Cada carpeta tiene una función específica dentro del sistema.



% --------------------------------------------------
% ARCHIVO PRINCIPAL
% --------------------------------------------------

\section{Archivo principal}

El archivo principal del proyecto es:

\begin{lstlisting}[caption={Archivo principal del documento},label={code:mainfile}]
main.tex
\end{lstlisting}

\source{Plantilla}

Este archivo controla la estructura general del documento
e incluye los distintos capítulos del trabajo.



% --------------------------------------------------
% FLUJO DE COMPILACIÓN
% --------------------------------------------------

\section{Flujo de compilación}

Durante la compilación del documento, \LaTeX{} sigue
una secuencia similar a la siguiente.

\begin{lstlisting}[caption={Flujo de compilación del documento},label={code:compileflow}]
main.tex

-> carga configuracion global
-> carga paquetes
-> carga estilos tipograficos
-> carga acronimos
-> genera frontmatter
-> compila capitulos
-> genera bibliografia
-> compila apendices
\end{lstlisting}

\source{Plantilla}



% --------------------------------------------------
% RESPONSABILIDAD DE LOS MODULOS
% --------------------------------------------------

\section{Responsabilidad de los módulos}

\begin{itemize}

\item \textbf{00\_config}  
Configuración global del documento.

\item \textbf{01\_frontmatter}  
Portada, dedicatoria, resúmenes.

\item \textbf{02\_chapters}  
Contenido principal.

\item \textbf{03\_figures}  
Imágenes del documento.

\item \textbf{04\_bibliography}  
Base de datos de referencias.

\item \textbf{05\_appendices}  
Apéndices del documento.

\item \textbf{06\_styles}  
Definición de estilos tipográficos.

\item \textbf{07\_acronyms}  
Definición de acrónimos.

\end{itemize}




% --------------------------------------------------
% BENEFICIOS DE LA ARQUITECTURA
% --------------------------------------------------

\section{Beneficios de la arquitectura}

Esta arquitectura modular ofrece varias ventajas.

\begin{itemize}

\item facilita el mantenimiento del proyecto
\item permite reutilizar la plantilla en múltiples documentos
\item separa claramente contenido y formato
\item reduce errores de compilación

\end{itemize}



\begin{examplebox}
Un autor puede escribir todos los capítulos del documento
sin necesidad de modificar los archivos de configuración
o los estilos tipográficos.
\end{examplebox}

%%%%%%%%%%%%%%%%%%%%%%%%%%%%%%%%%%%%%%%%%%%%%%%%%%%%%%%%%%
% FILE: chapter03_escribir_contenido.tex
%
% PURPOSE
% -------------------------------------------------------
% Explains how to write the main content of the document
% using the template structure.
%
%
% ARCHITECTURE ROLE
% -------------------------------------------------------
% User guide chapter.
%
% Describes the workflow authors should follow when
% writing documents with this template.
%
%
% DEPENDENCIES
% -------------------------------------------------------
% environments.tex
% listings
%
%
% NOTES FOR DEVELOPERS
% -------------------------------------------------------
% This chapter explains how users should interact with
% the template when writing content.
%
% It should not document template internals.
%
%%%%%%%%%%%%%%%%%%%%%%%%%%%%%%%%%%%%%%%%%%%%%%%%%%%%%%%%%%


\chapter{Cómo escribir contenido}

El contenido principal del documento se escribe en la
carpeta:

\texttt{02\_chapters}

Cada archivo dentro de esta carpeta representa un capítulo
independiente del documento.



% --------------------------------------------------
% ORGANIZACIÓN DE LOS CAPÍTULOS
% --------------------------------------------------

\section{Organización de los capítulos}

Un capítulo se define mediante el comando:

\begin{lstlisting}[caption={Creación de un capítulo},label={code:chapter}]
\chapter{Titulo del capitulo}
\end{lstlisting}

\source{Plantilla}

Cada capítulo debe almacenarse en un archivo independiente
dentro de la carpeta \texttt{02\_chapters}.



\begin{definitionbox}[Capítulo]
Un capítulo es la unidad principal de organización
del contenido dentro de un documento largo.
\end{definitionbox}



% --------------------------------------------------
% SECCIONES
% --------------------------------------------------

\section{Secciones}

Las secciones permiten dividir un capítulo en bloques
temáticos.

\begin{lstlisting}[caption={Creación de una sección},label={code:section}]
\section{Titulo de la seccion}
\end{lstlisting}

\source{Plantilla}



% --------------------------------------------------
% SUBSECCIONES
% --------------------------------------------------

\section{Subsecciones}

Las subsecciones permiten organizar el contenido con
mayor nivel de detalle.

\begin{lstlisting}[caption={Creación de una subsección},label={code:subsection}]
\subsection{Titulo de la subseccion}
\end{lstlisting}

\source{Plantilla}



\begin{notebox}
Se recomienda no utilizar más de tres niveles
de profundidad en la estructura del documento
para mantener la claridad del texto.
\end{notebox}



% --------------------------------------------------
% EJEMPLO DE ESTRUCTURA
% --------------------------------------------------

\section{Ejemplo de estructura}

Un capítulo típico puede organizarse de la siguiente forma.

\begin{lstlisting}[caption={Estructura típica de un capítulo},label={code:chapterstructure}]
\chapter{Introduccion}

\section{Contexto}

Texto introductorio.

\section{Objetivos}

\subsection{Objetivo general}

Descripcion.

\subsection{Objetivos especificos}

Descripcion.
\end{lstlisting}

\source{Plantilla}



% --------------------------------------------------
% FLUJO DE ESCRITURA
% --------------------------------------------------

\section{Flujo de escritura recomendado}

El flujo de trabajo recomendado al escribir un documento
con esta plantilla es el siguiente.

\begin{enumerate}

\item Crear un archivo nuevo en la carpeta
\texttt{02\_chapters}

\item Añadir el capítulo en \texttt{main.tex}

\item Escribir el contenido utilizando secciones
y subsecciones

\end{enumerate}



\begin{lstlisting}[caption={Carga de capítulos en main.tex},label={code:chapters}]
%%%%%%%%%%%%%%%%%%%%%%%%%%%%%%%%%%%%%%%%%%%%%%%%%%%%%%%%%%
% FILE: chapter01_introduccion.tex
%
% PURPOSE
% -------------------------------------------------------
% Introduces the LaTeX template and explains its
% objectives, philosophy and intended use.
%
%
% ARCHITECTURE ROLE
% -------------------------------------------------------
% First chapter of the template documentation.
%
% Provides conceptual context before explaining
% the technical structure of the project.
%
%
% DEPENDENCIES
% -------------------------------------------------------
% environments.tex
% (definitionbox, notebox, warningbox)
%
%
% NOTES FOR DEVELOPERS
% -------------------------------------------------------
% This chapter should remain conceptual.
%
% Do not include implementation details here.
% Technical explanations belong to later chapters.
%
%%%%%%%%%%%%%%%%%%%%%%%%%%%%%%%%%%%%%%%%%%%%%%%%%%%%%%%%%%


\chapter{Introducción}

Esta plantilla ha sido diseñada para facilitar la redacción de
documentos académicos y técnicos utilizando \LaTeX.

El objetivo principal es proporcionar una estructura modular
que permita separar claramente:

\begin{itemize}
\item configuración del documento
\item estilos tipográficos
\item contenido académico
\item recursos gráficos y bibliográficos
\end{itemize}

Gracias a esta separación, el autor puede concentrarse
exclusivamente en escribir el contenido del documento
sin preocuparse por el formato.



% --------------------------------------------------
% OBJETIVO DE LA PLANTILLA
% --------------------------------------------------

\section{Objetivo de la plantilla}

\begin{definitionbox}[Plantilla académica]
Una plantilla académica es una estructura predefinida
que automatiza la configuración tipográfica y
organizativa de un documento.
\end{definitionbox}

Esta plantilla está diseñada para:

\begin{itemize}
\item trabajos universitarios
\item informes técnicos
\item proyectos de ingeniería
\item Trabajos Fin de Grado (TFG)
\item Trabajos Fin de Máster (TFM)
\end{itemize}



% --------------------------------------------------
% FILOSOFÍA DE DISEÑO
% --------------------------------------------------

\section{Filosofía de diseño}

La plantilla sigue varios principios de diseño inspirados
en prácticas de ingeniería de software.

\begin{itemize}
\item modularidad
\item separación de responsabilidades
\item reutilización
\item documentación interna
\end{itemize}



\begin{notebox}
El objetivo no es únicamente generar un documento
correctamente formateado, sino también demostrar
buenas prácticas en la organización de proyectos
\LaTeX.
\end{notebox}



% --------------------------------------------------
% ORGANIZACIÓN DEL PROYECTO
% --------------------------------------------------

\section{Organización del proyecto}

El proyecto se organiza en varias carpetas que separan
la configuración del documento del contenido académico.



\begin{lstlisting}[caption={Estructura general del proyecto},label={code:estructura}]
main.tex

latex-university-template/
|-- main.tex
|-- 00_config/
|   |-- config.tex
|   `-- packages.tex
|-- 01_frontmatter/
|-- 02_chapters/
|-- 03_figures/
|-- 04_bibliography/
|-- 05_appendices/
|-- 06_styles/
`-- 07_acronyms/
\end{lstlisting}

\source{Plantilla}



% --------------------------------------------------
% USO RECOMENDADO
% --------------------------------------------------

\section{Uso recomendado}

El flujo de trabajo recomendado es el siguiente:

\begin{enumerate}
\item Configurar los metadatos del documento
\item Escribir el contenido en la carpeta \texttt{02\_chapters}
\item Añadir figuras en \texttt{03\_figures}
\item Añadir referencias en \texttt{05\_bibliography}
\end{enumerate}



\begin{warningbox}
Modificar directamente los archivos de configuración
o de estilos sin comprender su función puede producir
errores de compilación o inconsistencias tipográficas.
\end{warningbox}



% --------------------------------------------------
% METODOLOGÍA DE DOCUMENTACIÓN
% --------------------------------------------------

\section{Metodología de documentación}

Todos los archivos de la plantilla siguen una convención
de comentarios estructurados.

Cada archivo comienza con un bloque de documentación
que describe su función dentro del proyecto.

\begin{lstlisting}[caption={Bloque de documentación estándar},label={code:docblock}]
%%%%%%%%%%%%%%%%%%%%%%%%%%%%%%%%%%%%%%%%%%%%%%%%%%%%%%%%%%
% FILE: filename.tex
%
% PURPOSE
% Describe el proposito del archivo
%
% ARCHITECTURE ROLE
% Describe su funcion dentro del sistema
%
% DEPENDENCIES
% Archivos o paquetes necesarios
%
% NOTES FOR DEVELOPERS
% Indicaciones para mantenimiento
%%%%%%%%%%%%%%%%%%%%%%%%%%%%%%%%%%%%%%%%%%%%%%%%%%%%%%%%%%
\end{lstlisting}

\source{Plantilla}



\begin{examplebox}
Este sistema de documentación permite comprender
rápidamente la arquitectura del proyecto al revisar
los archivos de la plantilla.
\end{examplebox}

%%%%%%%%%%%%%%%%%%%%%%%%%%%%%%%%%%%%%%%%%%%%%%%%%%%%%%%%%%
% FILE: chapter02_arquitectura.tex
%
% PURPOSE
% -------------------------------------------------------
% Explains the internal architecture of the LaTeX template.
%
%
% ARCHITECTURE ROLE
% -------------------------------------------------------
% Technical documentation of the template structure.
%
%
% DEPENDENCIES
% -------------------------------------------------------
% environments.tex
%
%
% NOTES FOR DEVELOPERS
% -------------------------------------------------------
% If the directory structure of the template changes,
% this chapter must be updated accordingly.
%
%%%%%%%%%%%%%%%%%%%%%%%%%%%%%%%%%%%%%%%%%%%%%%%%%%%%%%%%%%


\chapter{Arquitectura del proyecto}

La plantilla está diseñada siguiendo principios de
arquitectura modular inspirados en prácticas de
ingeniería de software.

Cada componente del proyecto tiene una responsabilidad
claramente definida dentro del sistema.



% --------------------------------------------------
% PRINCIPIOS DE ARQUITECTURA
% --------------------------------------------------

\section{Principios de arquitectura}

\begin{definitionbox}[Arquitectura modular]
Una arquitectura modular divide un sistema en componentes
independientes que pueden mantenerse o modificarse de
forma aislada.
\end{definitionbox}

En esta plantilla, el proyecto se divide en varios módulos
que gestionan aspectos distintos del documento.



\begin{itemize}

\item configuración global
\item estilos tipográficos
\item contenido académico
\item recursos gráficos
\item bibliografía
\item acrónimos

\end{itemize}



\begin{notebox}
Separar estas responsabilidades permite modificar el
formato del documento sin alterar el contenido.
\end{notebox}



% --------------------------------------------------
% ESTRUCTURA DEL PROYECTO
% --------------------------------------------------

\section{Estructura del proyecto}

La estructura de carpetas del proyecto es la siguiente.

\begin{lstlisting}[caption={Estructura de directorios del proyecto},label={code:projectstructure}]
latex-university-template/
|-- main.tex
|-- 00_config/
|   |-- config.tex
|   `-- packages.tex
|-- 01_frontmatter/
|-- 02_chapters/
|-- 03_figures/
|-- 04_bibliography/
|-- 05_appendices/
|-- 06_styles/
`-- 07_acronyms/
\end{lstlisting}

\source{Plantilla}

Cada carpeta tiene una función específica dentro del sistema.



% --------------------------------------------------
% ARCHIVO PRINCIPAL
% --------------------------------------------------

\section{Archivo principal}

El archivo principal del proyecto es:

\begin{lstlisting}[caption={Archivo principal del documento},label={code:mainfile}]
main.tex
\end{lstlisting}

\source{Plantilla}

Este archivo controla la estructura general del documento
e incluye los distintos capítulos del trabajo.



% --------------------------------------------------
% FLUJO DE COMPILACIÓN
% --------------------------------------------------

\section{Flujo de compilación}

Durante la compilación del documento, \LaTeX{} sigue
una secuencia similar a la siguiente.

\begin{lstlisting}[caption={Flujo de compilación del documento},label={code:compileflow}]
main.tex

-> carga configuracion global
-> carga paquetes
-> carga estilos tipograficos
-> carga acronimos
-> genera frontmatter
-> compila capitulos
-> genera bibliografia
-> compila apendices
\end{lstlisting}

\source{Plantilla}



% --------------------------------------------------
% RESPONSABILIDAD DE LOS MODULOS
% --------------------------------------------------

\section{Responsabilidad de los módulos}

\begin{itemize}

\item \textbf{00\_config}  
Configuración global del documento.

\item \textbf{01\_frontmatter}  
Portada, dedicatoria, resúmenes.

\item \textbf{02\_chapters}  
Contenido principal.

\item \textbf{03\_figures}  
Imágenes del documento.

\item \textbf{04\_bibliography}  
Base de datos de referencias.

\item \textbf{05\_appendices}  
Apéndices del documento.

\item \textbf{06\_styles}  
Definición de estilos tipográficos.

\item \textbf{07\_acronyms}  
Definición de acrónimos.

\end{itemize}




% --------------------------------------------------
% BENEFICIOS DE LA ARQUITECTURA
% --------------------------------------------------

\section{Beneficios de la arquitectura}

Esta arquitectura modular ofrece varias ventajas.

\begin{itemize}

\item facilita el mantenimiento del proyecto
\item permite reutilizar la plantilla en múltiples documentos
\item separa claramente contenido y formato
\item reduce errores de compilación

\end{itemize}



\begin{examplebox}
Un autor puede escribir todos los capítulos del documento
sin necesidad de modificar los archivos de configuración
o los estilos tipográficos.
\end{examplebox}

%%%%%%%%%%%%%%%%%%%%%%%%%%%%%%%%%%%%%%%%%%%%%%%%%%%%%%%%%%
% FILE: chapter03_escribir_contenido.tex
%
% PURPOSE
% -------------------------------------------------------
% Explains how to write the main content of the document
% using the template structure.
%
%
% ARCHITECTURE ROLE
% -------------------------------------------------------
% User guide chapter.
%
% Describes the workflow authors should follow when
% writing documents with this template.
%
%
% DEPENDENCIES
% -------------------------------------------------------
% environments.tex
% listings
%
%
% NOTES FOR DEVELOPERS
% -------------------------------------------------------
% This chapter explains how users should interact with
% the template when writing content.
%
% It should not document template internals.
%
%%%%%%%%%%%%%%%%%%%%%%%%%%%%%%%%%%%%%%%%%%%%%%%%%%%%%%%%%%


\chapter{Cómo escribir contenido}

El contenido principal del documento se escribe en la
carpeta:

\texttt{02\_chapters}

Cada archivo dentro de esta carpeta representa un capítulo
independiente del documento.



% --------------------------------------------------
% ORGANIZACIÓN DE LOS CAPÍTULOS
% --------------------------------------------------

\section{Organización de los capítulos}

Un capítulo se define mediante el comando:

\begin{lstlisting}[caption={Creación de un capítulo},label={code:chapter}]
\chapter{Titulo del capitulo}
\end{lstlisting}

\source{Plantilla}

Cada capítulo debe almacenarse en un archivo independiente
dentro de la carpeta \texttt{02\_chapters}.



\begin{definitionbox}[Capítulo]
Un capítulo es la unidad principal de organización
del contenido dentro de un documento largo.
\end{definitionbox}



% --------------------------------------------------
% SECCIONES
% --------------------------------------------------

\section{Secciones}

Las secciones permiten dividir un capítulo en bloques
temáticos.

\begin{lstlisting}[caption={Creación de una sección},label={code:section}]
\section{Titulo de la seccion}
\end{lstlisting}

\source{Plantilla}



% --------------------------------------------------
% SUBSECCIONES
% --------------------------------------------------

\section{Subsecciones}

Las subsecciones permiten organizar el contenido con
mayor nivel de detalle.

\begin{lstlisting}[caption={Creación de una subsección},label={code:subsection}]
\subsection{Titulo de la subseccion}
\end{lstlisting}

\source{Plantilla}



\begin{notebox}
Se recomienda no utilizar más de tres niveles
de profundidad en la estructura del documento
para mantener la claridad del texto.
\end{notebox}



% --------------------------------------------------
% EJEMPLO DE ESTRUCTURA
% --------------------------------------------------

\section{Ejemplo de estructura}

Un capítulo típico puede organizarse de la siguiente forma.

\begin{lstlisting}[caption={Estructura típica de un capítulo},label={code:chapterstructure}]
\chapter{Introduccion}

\section{Contexto}

Texto introductorio.

\section{Objetivos}

\subsection{Objetivo general}

Descripcion.

\subsection{Objetivos especificos}

Descripcion.
\end{lstlisting}

\source{Plantilla}



% --------------------------------------------------
% FLUJO DE ESCRITURA
% --------------------------------------------------

\section{Flujo de escritura recomendado}

El flujo de trabajo recomendado al escribir un documento
con esta plantilla es el siguiente.

\begin{enumerate}

\item Crear un archivo nuevo en la carpeta
\texttt{02\_chapters}

\item Añadir el capítulo en \texttt{main.tex}

\item Escribir el contenido utilizando secciones
y subsecciones

\end{enumerate}



\begin{lstlisting}[caption={Carga de capítulos en main.tex},label={code:chapters}]
%%%%%%%%%%%%%%%%%%%%%%%%%%%%%%%%%%%%%%%%%%%%%%%%%%%%%%%%%%
% FILE: chapter01_introduccion.tex
%
% PURPOSE
% -------------------------------------------------------
% Introduces the LaTeX template and explains its
% objectives, philosophy and intended use.
%
%
% ARCHITECTURE ROLE
% -------------------------------------------------------
% First chapter of the template documentation.
%
% Provides conceptual context before explaining
% the technical structure of the project.
%
%
% DEPENDENCIES
% -------------------------------------------------------
% environments.tex
% (definitionbox, notebox, warningbox)
%
%
% NOTES FOR DEVELOPERS
% -------------------------------------------------------
% This chapter should remain conceptual.
%
% Do not include implementation details here.
% Technical explanations belong to later chapters.
%
%%%%%%%%%%%%%%%%%%%%%%%%%%%%%%%%%%%%%%%%%%%%%%%%%%%%%%%%%%


\chapter{Introducción}

Esta plantilla ha sido diseñada para facilitar la redacción de
documentos académicos y técnicos utilizando \LaTeX.

El objetivo principal es proporcionar una estructura modular
que permita separar claramente:

\begin{itemize}
\item configuración del documento
\item estilos tipográficos
\item contenido académico
\item recursos gráficos y bibliográficos
\end{itemize}

Gracias a esta separación, el autor puede concentrarse
exclusivamente en escribir el contenido del documento
sin preocuparse por el formato.



% --------------------------------------------------
% OBJETIVO DE LA PLANTILLA
% --------------------------------------------------

\section{Objetivo de la plantilla}

\begin{definitionbox}[Plantilla académica]
Una plantilla académica es una estructura predefinida
que automatiza la configuración tipográfica y
organizativa de un documento.
\end{definitionbox}

Esta plantilla está diseñada para:

\begin{itemize}
\item trabajos universitarios
\item informes técnicos
\item proyectos de ingeniería
\item Trabajos Fin de Grado (TFG)
\item Trabajos Fin de Máster (TFM)
\end{itemize}



% --------------------------------------------------
% FILOSOFÍA DE DISEÑO
% --------------------------------------------------

\section{Filosofía de diseño}

La plantilla sigue varios principios de diseño inspirados
en prácticas de ingeniería de software.

\begin{itemize}
\item modularidad
\item separación de responsabilidades
\item reutilización
\item documentación interna
\end{itemize}



\begin{notebox}
El objetivo no es únicamente generar un documento
correctamente formateado, sino también demostrar
buenas prácticas en la organización de proyectos
\LaTeX.
\end{notebox}



% --------------------------------------------------
% ORGANIZACIÓN DEL PROYECTO
% --------------------------------------------------

\section{Organización del proyecto}

El proyecto se organiza en varias carpetas que separan
la configuración del documento del contenido académico.



\begin{lstlisting}[caption={Estructura general del proyecto},label={code:estructura}]
main.tex

latex-university-template/
|-- main.tex
|-- 00_config/
|   |-- config.tex
|   `-- packages.tex
|-- 01_frontmatter/
|-- 02_chapters/
|-- 03_figures/
|-- 04_bibliography/
|-- 05_appendices/
|-- 06_styles/
`-- 07_acronyms/
\end{lstlisting}

\source{Plantilla}



% --------------------------------------------------
% USO RECOMENDADO
% --------------------------------------------------

\section{Uso recomendado}

El flujo de trabajo recomendado es el siguiente:

\begin{enumerate}
\item Configurar los metadatos del documento
\item Escribir el contenido en la carpeta \texttt{02\_chapters}
\item Añadir figuras en \texttt{03\_figures}
\item Añadir referencias en \texttt{05\_bibliography}
\end{enumerate}



\begin{warningbox}
Modificar directamente los archivos de configuración
o de estilos sin comprender su función puede producir
errores de compilación o inconsistencias tipográficas.
\end{warningbox}



% --------------------------------------------------
% METODOLOGÍA DE DOCUMENTACIÓN
% --------------------------------------------------

\section{Metodología de documentación}

Todos los archivos de la plantilla siguen una convención
de comentarios estructurados.

Cada archivo comienza con un bloque de documentación
que describe su función dentro del proyecto.

\begin{lstlisting}[caption={Bloque de documentación estándar},label={code:docblock}]
%%%%%%%%%%%%%%%%%%%%%%%%%%%%%%%%%%%%%%%%%%%%%%%%%%%%%%%%%%
% FILE: filename.tex
%
% PURPOSE
% Describe el proposito del archivo
%
% ARCHITECTURE ROLE
% Describe su funcion dentro del sistema
%
% DEPENDENCIES
% Archivos o paquetes necesarios
%
% NOTES FOR DEVELOPERS
% Indicaciones para mantenimiento
%%%%%%%%%%%%%%%%%%%%%%%%%%%%%%%%%%%%%%%%%%%%%%%%%%%%%%%%%%
\end{lstlisting}

\source{Plantilla}



\begin{examplebox}
Este sistema de documentación permite comprender
rápidamente la arquitectura del proyecto al revisar
los archivos de la plantilla.
\end{examplebox}

%%%%%%%%%%%%%%%%%%%%%%%%%%%%%%%%%%%%%%%%%%%%%%%%%%%%%%%%%%
% FILE: chapter02_arquitectura.tex
%
% PURPOSE
% -------------------------------------------------------
% Explains the internal architecture of the LaTeX template.
%
%
% ARCHITECTURE ROLE
% -------------------------------------------------------
% Technical documentation of the template structure.
%
%
% DEPENDENCIES
% -------------------------------------------------------
% environments.tex
%
%
% NOTES FOR DEVELOPERS
% -------------------------------------------------------
% If the directory structure of the template changes,
% this chapter must be updated accordingly.
%
%%%%%%%%%%%%%%%%%%%%%%%%%%%%%%%%%%%%%%%%%%%%%%%%%%%%%%%%%%


\chapter{Arquitectura del proyecto}

La plantilla está diseñada siguiendo principios de
arquitectura modular inspirados en prácticas de
ingeniería de software.

Cada componente del proyecto tiene una responsabilidad
claramente definida dentro del sistema.



% --------------------------------------------------
% PRINCIPIOS DE ARQUITECTURA
% --------------------------------------------------

\section{Principios de arquitectura}

\begin{definitionbox}[Arquitectura modular]
Una arquitectura modular divide un sistema en componentes
independientes que pueden mantenerse o modificarse de
forma aislada.
\end{definitionbox}

En esta plantilla, el proyecto se divide en varios módulos
que gestionan aspectos distintos del documento.



\begin{itemize}

\item configuración global
\item estilos tipográficos
\item contenido académico
\item recursos gráficos
\item bibliografía
\item acrónimos

\end{itemize}



\begin{notebox}
Separar estas responsabilidades permite modificar el
formato del documento sin alterar el contenido.
\end{notebox}



% --------------------------------------------------
% ESTRUCTURA DEL PROYECTO
% --------------------------------------------------

\section{Estructura del proyecto}

La estructura de carpetas del proyecto es la siguiente.

\begin{lstlisting}[caption={Estructura de directorios del proyecto},label={code:projectstructure}]
latex-university-template/
|-- main.tex
|-- 00_config/
|   |-- config.tex
|   `-- packages.tex
|-- 01_frontmatter/
|-- 02_chapters/
|-- 03_figures/
|-- 04_bibliography/
|-- 05_appendices/
|-- 06_styles/
`-- 07_acronyms/
\end{lstlisting}

\source{Plantilla}

Cada carpeta tiene una función específica dentro del sistema.



% --------------------------------------------------
% ARCHIVO PRINCIPAL
% --------------------------------------------------

\section{Archivo principal}

El archivo principal del proyecto es:

\begin{lstlisting}[caption={Archivo principal del documento},label={code:mainfile}]
main.tex
\end{lstlisting}

\source{Plantilla}

Este archivo controla la estructura general del documento
e incluye los distintos capítulos del trabajo.



% --------------------------------------------------
% FLUJO DE COMPILACIÓN
% --------------------------------------------------

\section{Flujo de compilación}

Durante la compilación del documento, \LaTeX{} sigue
una secuencia similar a la siguiente.

\begin{lstlisting}[caption={Flujo de compilación del documento},label={code:compileflow}]
main.tex

-> carga configuracion global
-> carga paquetes
-> carga estilos tipograficos
-> carga acronimos
-> genera frontmatter
-> compila capitulos
-> genera bibliografia
-> compila apendices
\end{lstlisting}

\source{Plantilla}



% --------------------------------------------------
% RESPONSABILIDAD DE LOS MODULOS
% --------------------------------------------------

\section{Responsabilidad de los módulos}

\begin{itemize}

\item \textbf{00\_config}  
Configuración global del documento.

\item \textbf{01\_frontmatter}  
Portada, dedicatoria, resúmenes.

\item \textbf{02\_chapters}  
Contenido principal.

\item \textbf{03\_figures}  
Imágenes del documento.

\item \textbf{04\_bibliography}  
Base de datos de referencias.

\item \textbf{05\_appendices}  
Apéndices del documento.

\item \textbf{06\_styles}  
Definición de estilos tipográficos.

\item \textbf{07\_acronyms}  
Definición de acrónimos.

\end{itemize}




% --------------------------------------------------
% BENEFICIOS DE LA ARQUITECTURA
% --------------------------------------------------

\section{Beneficios de la arquitectura}

Esta arquitectura modular ofrece varias ventajas.

\begin{itemize}

\item facilita el mantenimiento del proyecto
\item permite reutilizar la plantilla en múltiples documentos
\item separa claramente contenido y formato
\item reduce errores de compilación

\end{itemize}



\begin{examplebox}
Un autor puede escribir todos los capítulos del documento
sin necesidad de modificar los archivos de configuración
o los estilos tipográficos.
\end{examplebox}

%%%%%%%%%%%%%%%%%%%%%%%%%%%%%%%%%%%%%%%%%%%%%%%%%%%%%%%%%%
% FILE: chapter03_escribir_contenido.tex
%
% PURPOSE
% -------------------------------------------------------
% Explains how to write the main content of the document
% using the template structure.
%
%
% ARCHITECTURE ROLE
% -------------------------------------------------------
% User guide chapter.
%
% Describes the workflow authors should follow when
% writing documents with this template.
%
%
% DEPENDENCIES
% -------------------------------------------------------
% environments.tex
% listings
%
%
% NOTES FOR DEVELOPERS
% -------------------------------------------------------
% This chapter explains how users should interact with
% the template when writing content.
%
% It should not document template internals.
%
%%%%%%%%%%%%%%%%%%%%%%%%%%%%%%%%%%%%%%%%%%%%%%%%%%%%%%%%%%


\chapter{Cómo escribir contenido}

El contenido principal del documento se escribe en la
carpeta:

\texttt{02\_chapters}

Cada archivo dentro de esta carpeta representa un capítulo
independiente del documento.



% --------------------------------------------------
% ORGANIZACIÓN DE LOS CAPÍTULOS
% --------------------------------------------------

\section{Organización de los capítulos}

Un capítulo se define mediante el comando:

\begin{lstlisting}[caption={Creación de un capítulo},label={code:chapter}]
\chapter{Titulo del capitulo}
\end{lstlisting}

\source{Plantilla}

Cada capítulo debe almacenarse en un archivo independiente
dentro de la carpeta \texttt{02\_chapters}.



\begin{definitionbox}[Capítulo]
Un capítulo es la unidad principal de organización
del contenido dentro de un documento largo.
\end{definitionbox}



% --------------------------------------------------
% SECCIONES
% --------------------------------------------------

\section{Secciones}

Las secciones permiten dividir un capítulo en bloques
temáticos.

\begin{lstlisting}[caption={Creación de una sección},label={code:section}]
\section{Titulo de la seccion}
\end{lstlisting}

\source{Plantilla}



% --------------------------------------------------
% SUBSECCIONES
% --------------------------------------------------

\section{Subsecciones}

Las subsecciones permiten organizar el contenido con
mayor nivel de detalle.

\begin{lstlisting}[caption={Creación de una subsección},label={code:subsection}]
\subsection{Titulo de la subseccion}
\end{lstlisting}

\source{Plantilla}



\begin{notebox}
Se recomienda no utilizar más de tres niveles
de profundidad en la estructura del documento
para mantener la claridad del texto.
\end{notebox}



% --------------------------------------------------
% EJEMPLO DE ESTRUCTURA
% --------------------------------------------------

\section{Ejemplo de estructura}

Un capítulo típico puede organizarse de la siguiente forma.

\begin{lstlisting}[caption={Estructura típica de un capítulo},label={code:chapterstructure}]
\chapter{Introduccion}

\section{Contexto}

Texto introductorio.

\section{Objetivos}

\subsection{Objetivo general}

Descripcion.

\subsection{Objetivos especificos}

Descripcion.
\end{lstlisting}

\source{Plantilla}



% --------------------------------------------------
% FLUJO DE ESCRITURA
% --------------------------------------------------

\section{Flujo de escritura recomendado}

El flujo de trabajo recomendado al escribir un documento
con esta plantilla es el siguiente.

\begin{enumerate}

\item Crear un archivo nuevo en la carpeta
\texttt{02\_chapters}

\item Añadir el capítulo en \texttt{main.tex}

\item Escribir el contenido utilizando secciones
y subsecciones

\end{enumerate}



\begin{lstlisting}[caption={Carga de capítulos en main.tex},label={code:chapters}]
\input{02_chapters/chapter01}
\input{02_chapters/chapter02}
\input{02_chapters/chapter03}
\end{lstlisting}

\source{Plantilla}



% --------------------------------------------------
% BUENAS PRÁCTICAS DE ESCRITURA
% --------------------------------------------------

\section{Buenas prácticas de escritura}

Para mantener la coherencia del documento se recomienda
seguir estas prácticas.

\begin{itemize}

\item cada capítulo en un archivo independiente
\item nombres de archivos descriptivos
\item evitar capítulos excesivamente largos
\item mantener una jerarquía clara de secciones

\end{itemize}



\begin{warningbox}
Escribir todo el contenido directamente en
\texttt{main.tex} rompe la arquitectura modular
de la plantilla y dificulta el mantenimiento
del documento.
\end{warningbox}



% --------------------------------------------------
% EJEMPLO DE USO
% --------------------------------------------------

\section{Ejemplo}

\begin{examplebox}
Un documento típico puede contener capítulos como:

Introducción

Marco teórico

Metodología

Resultados

Conclusiones
\end{examplebox}

\end{lstlisting}

\source{Plantilla}



% --------------------------------------------------
% BUENAS PRÁCTICAS DE ESCRITURA
% --------------------------------------------------

\section{Buenas prácticas de escritura}

Para mantener la coherencia del documento se recomienda
seguir estas prácticas.

\begin{itemize}

\item cada capítulo en un archivo independiente
\item nombres de archivos descriptivos
\item evitar capítulos excesivamente largos
\item mantener una jerarquía clara de secciones

\end{itemize}



\begin{warningbox}
Escribir todo el contenido directamente en
\texttt{main.tex} rompe la arquitectura modular
de la plantilla y dificulta el mantenimiento
del documento.
\end{warningbox}



% --------------------------------------------------
% EJEMPLO DE USO
% --------------------------------------------------

\section{Ejemplo}

\begin{examplebox}
Un documento típico puede contener capítulos como:

Introducción

Marco teórico

Metodología

Resultados

Conclusiones
\end{examplebox}

\end{lstlisting}

\source{Plantilla}



% --------------------------------------------------
% BUENAS PRÁCTICAS DE ESCRITURA
% --------------------------------------------------

\section{Buenas prácticas de escritura}

Para mantener la coherencia del documento se recomienda
seguir estas prácticas.

\begin{itemize}

\item cada capítulo en un archivo independiente
\item nombres de archivos descriptivos
\item evitar capítulos excesivamente largos
\item mantener una jerarquía clara de secciones

\end{itemize}



\begin{warningbox}
Escribir todo el contenido directamente en
\texttt{main.tex} rompe la arquitectura modular
de la plantilla y dificulta el mantenimiento
del documento.
\end{warningbox}



% --------------------------------------------------
% EJEMPLO DE USO
% --------------------------------------------------

\section{Ejemplo}

\begin{examplebox}
Un documento típico puede contener capítulos como:

Introducción

Marco teórico

Metodología

Resultados

Conclusiones
\end{examplebox}

\end{lstlisting}

\source{Plantilla}



% --------------------------------------------------
% BUENAS PRÁCTICAS DE ESCRITURA
% --------------------------------------------------

\section{Buenas prácticas de escritura}

Para mantener la coherencia del documento se recomienda
seguir estas prácticas.

\begin{itemize}

\item cada capítulo en un archivo independiente
\item nombres de archivos descriptivos
\item evitar capítulos excesivamente largos
\item mantener una jerarquía clara de secciones

\end{itemize}



\begin{warningbox}
Escribir todo el contenido directamente en
\texttt{main.tex} rompe la arquitectura modular
de la plantilla y dificulta el mantenimiento
del documento.
\end{warningbox}



% --------------------------------------------------
% EJEMPLO DE USO
% --------------------------------------------------

\section{Ejemplo}

\begin{examplebox}
Un documento típico puede contener capítulos como:

Introducción

Marco teórico

Metodología

Resultados

Conclusiones
\end{examplebox}
